\documentclass[12pt,a4paper]{article}
\usepackage[utf8]{vietnam}
\usepackage{amsmath,amssymb}
\usepackage{tabularx}
\usepackage{tikz}
\usepackage{xcolor}
\usepackage[top=0.6cm, bottom=0.6cm, left=0.6cm, right=0.6cm]{geometry}
\usepackage[dethi]{ex_test}

\setlength{\parskip}{0pt}

% Dòng kẻ chấm mờ đều đặn chuẩn viết tay (khoảng cách 7.2mm rộng rãi)
\newcommand{\dongke}[1]{%
  \begingroup%
  \setlength{\parskip}{0pt}%
  \baselineskip=20.5pt%
  \vspace{1.5mm}%
  \foreach \i in {1,...,#1}{%
    \noindent{\color{gray!60}\dotfill}\par%
  }%
  \vspace{1mm}%
  \endgroup%
}

% Định nghĩa ô bài tập tự luận trong bảng
\newcommand{\oBaiTap}[3][6]{%
  \vspace{1.5mm}%
  \noindent\textbf{#2} #3\par%
  \dongke{#1}%
}

% Câu trắc nghiệm 1 cột: 4 đáp án dàn hàng ngang trên 1 dòng
\newcommand{\cauTN}[6]{% \cauTN{số câu}{đề}{A}{B}{C}{D}
  \noindent\textbf{Câu #1.} #2\par\vspace{0.4mm}%
  \noindent%
  \makebox[0.25\linewidth][l]{\textbf{A.} #3}%
  \makebox[0.25\linewidth][l]{\textbf{B.} #4}%
  \makebox[0.25\linewidth][l]{\textbf{C.} #5}%
  \makebox[0.25\linewidth][l]{\textbf{D.} #6}\par\vspace{0.8mm}%
}

\begin{document}
\pagestyle{empty}
% =============================================================================
% ========================= TOÁN 6 - TUẦN 26 - BUỔI 02: MÃ ĐỀ 101 =============
% ========================= MẶT TRƯỚC (TRANG 1/2) ==============================
% =============================================================================

\noindent
\begin{tabularx}{\linewidth}{|X|c|}
\hline
\textbf{PHIẾU KIỂM TRA ĐỊNH KỲ TOÁN 6 - TUẦN 26 - BUỔI 02} & \textbf{Mã đề: 101 - Trang 1/2} \\
\textit{Chủ đề: Bài 35: Trung điểm của đoạn thẳng} & \textit{Thời gian: 45 phút} \\
\hline
\end{tabularx}

\vspace{1mm}

\noindent
\begin{tabularx}{\linewidth}{|X|}
\hline
\rule[-3.5mm]{0pt}{10mm}\textbf{Họ và tên học sinh:} \hfill \textbf{Mã HS:} \makebox[3.5cm]{} \\
\hline
\end{tabularx}

\vspace{1mm}

\noindent
\begin{tabularx}{\linewidth}{|p{3.2cm}|X|}
\hline
\centering\textbf{ĐIỂM SỐ} & \centering\textbf{LỜI NHẬN XÉT CỦA THẦY CÔ} \tabularnewline
\hline
\rule{0pt}{1.3cm} & \tabularnewline
\hline
\end{tabularx}

\vspace{1.2mm}

\noindent\textbf{BẢNG ĐIỀN ĐÁP ÁN TRẮC NGHIỆM (Mỗi câu đúng 0,5 điểm):}
\vspace{0.8mm}

\noindent
\begingroup
\setlength{\tabcolsep}{3pt}
\renewcommand{\arraystretch}{1.35}
\begin{tabularx}{\linewidth}{|c|*{5}{>{\centering\arraybackslash}X|}}
\hline
\textbf{Câu} & \textbf{1} & \textbf{2} & \textbf{3} & \textbf{4} & \textbf{5} \\
\hline
\textbf{Đáp án} & & & & & \\
\hline
\end{tabularx}
\endgroup

\vspace{1.5mm}

\noindent\textbf{\large PHẦN I. TRẮC NGHIỆM KHÁCH QUAN (2,5 điểm)}
\vspace{1mm}

\cauTN{1}{Nội dung câu hỏi trắc nghiệm số 1 đang được biên soạn...}{Phương án A}{Phương án B}{Phương án C}{Phương án D}

\cauTN{2}{Nội dung câu hỏi trắc nghiệm số 2 đang được biên soạn...}{Phương án A}{Phương án B}{Phương án C}{Phương án D}

\cauTN{3}{Nội dung câu hỏi trắc nghiệm số 3 đang được biên soạn...}{Phương án A}{Phương án B}{Phương án C}{Phương án D}

\cauTN{4}{Nội dung câu hỏi trắc nghiệm số 4 đang được biên soạn...}{Phương án A}{Phương án B}{Phương án C}{Phương án D}

\cauTN{5}{Nội dung câu hỏi trắc nghiệm số 5 đang được biên soạn...}{Phương án A}{Phương án B}{Phương án C}{Phương án D}

\vspace{1.2mm}

\noindent\textbf{\large PHẦN II. TỰ LUẬN (7,5 điểm)}

\noindent\textbf{Câu 1 (2,5 điểm). Thực hiện phép tính (hợp lý nếu có thể):}
\vspace{0.8mm}

\noindent
\begin{tabularx}{\linewidth}{|X|X|}
\hline
\oBaiTap[6]{a)}{Nội dung ý a đang được cập nhật} & 
\oBaiTap[6]{b)}{Nội dung ý b đang được cập nhật} \\
\hline
\oBaiTap[6]{c)}{Nội dung ý c đang được cập nhật} & 
\oBaiTap[6]{d)}{Nội dung ý d đang được cập nhật} \\
\hline
\end{tabularx}

\newpage
% =============================================================================
% ========================= TOÁN 6 - TUẦN 26 - BUỔI 02: MÃ ĐỀ 101 =============
% ========================= MẶT SAU (TRANG 2/2) ================================
% =============================================================================

\noindent
\begin{tabularx}{\linewidth}{|X|c|}
\hline
\textbf{PHẦN II. TỰ LUẬN (Tiếp theo) - TOÁN 6 TUẦN 26 BUỔI 02} & \textbf{Mã đề: 101 - Trang 2/2} \\
\hline
\end{tabularx}

\vspace{2.5mm}

\noindent\textbf{\large Câu 2 (2,5 điểm). Tìm số tự nhiên $x$, biết:}
\vspace{0.8mm}

\noindent
\begin{tabularx}{\linewidth}{|X|X|}
\hline
\oBaiTap[6]{a)}{Nội dung ý a đang được cập nhật} & 
\oBaiTap[6]{b)}{Nội dung ý b đang được cập nhật} \\
\hline
\oBaiTap[6]{c)}{Nội dung ý c đang được cập nhật} & 
\oBaiTap[6]{d)}{Nội dung ý d đang được cập nhật} \\
\hline
\end{tabularx}

\vspace{2.5mm}

\noindent\textbf{\large Câu 3 (2,5 điểm). Bài toán thực tế:}
\vspace{0.8mm}

\noindent
\begin{tabularx}{\linewidth}{|X|}
\hline
\vspace{0.5mm}
Nội dung bài toán thực tế đang được cập nhật. Học sinh tóm tắt đề và trình bày bài giải chi tiết bên dưới.\par
\textbf{a) (1,25 điểm)} Câu hỏi ý a...\par
\textbf{b) (1,25 điểm)} Câu hỏi ý b...\par
\vspace{0.4mm}
\noindent\textit{Bài giải:}\dongke{14}
\\ \hline
\end{tabularx}

\newpage
% =============================================================================
% ========================= TOÁN 6 - TUẦN 26 - BUỔI 02: MÃ ĐỀ 102 =============
% ========================= MẶT TRƯỚC (TRANG 1/2) ==============================
% =============================================================================

\noindent
\begin{tabularx}{\linewidth}{|X|c|}
\hline
\textbf{PHIẾU KIỂM TRA ĐỊNH KỲ TOÁN 6 - TUẦN 26 - BUỔI 02} & \textbf{Mã đề: 102 - Trang 1/2} \\
\textit{Chủ đề: Bài 35: Trung điểm của đoạn thẳng} & \textit{Thời gian: 45 phút} \\
\hline
\end{tabularx}

\vspace{1mm}

\noindent
\begin{tabularx}{\linewidth}{|X|}
\hline
\rule[-3.5mm]{0pt}{10mm}\textbf{Họ và tên học sinh:} \hfill \textbf{Mã HS:} \makebox[3.5cm]{} \\
\hline
\end{tabularx}

\vspace{1mm}

\noindent
\begin{tabularx}{\linewidth}{|p{3.2cm}|X|}
\hline
\centering\textbf{ĐIỂM SỐ} & \centering\textbf{LỜI NHẬN XÉT CỦA THẦY CÔ} \tabularnewline
\hline
\rule{0pt}{1.3cm} & \tabularnewline
\hline
\end{tabularx}

\vspace{1.2mm}

\noindent\textbf{BẢNG ĐIỀN ĐÁP ÁN TRẮC NGHIỆM (Mỗi câu đúng 0,5 điểm):}
\vspace{0.8mm}

\noindent
\begingroup
\setlength{\tabcolsep}{3pt}
\renewcommand{\arraystretch}{1.35}
\begin{tabularx}{\linewidth}{|c|*{5}{>{\centering\arraybackslash}X|}}
\hline
\textbf{Câu} & \textbf{1} & \textbf{2} & \textbf{3} & \textbf{4} & \textbf{5} \\
\hline
\textbf{Đáp án} & & & & & \\
\hline
\end{tabularx}
\endgroup

\vspace{1.5mm}

\noindent\textbf{\large PHẦN I. TRẮC NGHIỆM KHÁCH QUAN (2,5 điểm)}
\vspace{1mm}

\cauTN{1}{Nội dung câu hỏi trắc nghiệm số 1 đang được biên soạn...}{Phương án A}{Phương án B}{Phương án C}{Phương án D}

\cauTN{2}{Nội dung câu hỏi trắc nghiệm số 2 đang được biên soạn...}{Phương án A}{Phương án B}{Phương án C}{Phương án D}

\cauTN{3}{Nội dung câu hỏi trắc nghiệm số 3 đang được biên soạn...}{Phương án A}{Phương án B}{Phương án C}{Phương án D}

\cauTN{4}{Nội dung câu hỏi trắc nghiệm số 4 đang được biên soạn...}{Phương án A}{Phương án B}{Phương án C}{Phương án D}

\cauTN{5}{Nội dung câu hỏi trắc nghiệm số 5 đang được biên soạn...}{Phương án A}{Phương án B}{Phương án C}{Phương án D}

\vspace{1.2mm}

\noindent\textbf{\large PHẦN II. TỰ LUẬN (7,5 điểm)}

\noindent\textbf{Câu 1 (2,5 điểm). Thực hiện phép tính (hợp lý nếu có thể):}
\vspace{0.8mm}

\noindent
\begin{tabularx}{\linewidth}{|X|X|}
\hline
\oBaiTap[6]{a)}{Nội dung ý a đang được cập nhật} & 
\oBaiTap[6]{b)}{Nội dung ý b đang được cập nhật} \\
\hline
\oBaiTap[6]{c)}{Nội dung ý c đang được cập nhật} & 
\oBaiTap[6]{d)}{Nội dung ý d đang được cập nhật} \\
\hline
\end{tabularx}

\newpage
% =============================================================================
% ========================= TOÁN 6 - TUẦN 26 - BUỔI 02: MÃ ĐỀ 102 =============
% ========================= MẶT SAU (TRANG 2/2) ================================
% =============================================================================

\noindent
\begin{tabularx}{\linewidth}{|X|c|}
\hline
\textbf{PHẦN II. TỰ LUẬN (Tiếp theo) - TOÁN 6 TUẦN 26 BUỔI 02} & \textbf{Mã đề: 102 - Trang 2/2} \\
\hline
\end{tabularx}

\vspace{2.5mm}

\noindent\textbf{\large Câu 2 (2,5 điểm). Tìm số tự nhiên $x$, biết:}
\vspace{0.8mm}

\noindent
\begin{tabularx}{\linewidth}{|X|X|}
\hline
\oBaiTap[6]{a)}{Nội dung ý a đang được cập nhật} & 
\oBaiTap[6]{b)}{Nội dung ý b đang được cập nhật} \\
\hline
\oBaiTap[6]{c)}{Nội dung ý c đang được cập nhật} & 
\oBaiTap[6]{d)}{Nội dung ý d đang được cập nhật} \\
\hline
\end{tabularx}

\vspace{2.5mm}

\noindent\textbf{\large Câu 3 (2,5 điểm). Bài toán thực tế:}
\vspace{0.8mm}

\noindent
\begin{tabularx}{\linewidth}{|X|}
\hline
\vspace{0.5mm}
Nội dung bài toán thực tế đang được cập nhật. Học sinh tóm tắt đề và trình bày bài giải chi tiết bên dưới.\par
\textbf{a) (1,25 điểm)} Câu hỏi ý a...\par
\textbf{b) (1,25 điểm)} Câu hỏi ý b...\par
\vspace{0.4mm}
\noindent\textit{Bài giải:}\dongke{14}
\\ \hline
\end{tabularx}

\newpage
% =============================================================================
% ========================= TOÁN 6 - TUẦN 26 - BUỔI 02: MÃ ĐỀ 103 =============
% ========================= MẶT TRƯỚC (TRANG 1/2) ==============================
% =============================================================================

\noindent
\begin{tabularx}{\linewidth}{|X|c|}
\hline
\textbf{PHIẾU KIỂM TRA ĐỊNH KỲ TOÁN 6 - TUẦN 26 - BUỔI 02} & \textbf{Mã đề: 103 - Trang 1/2} \\
\textit{Chủ đề: Bài 35: Trung điểm của đoạn thẳng} & \textit{Thời gian: 45 phút} \\
\hline
\end{tabularx}

\vspace{1mm}

\noindent
\begin{tabularx}{\linewidth}{|X|}
\hline
\rule[-3.5mm]{0pt}{10mm}\textbf{Họ và tên học sinh:} \hfill \textbf{Mã HS:} \makebox[3.5cm]{} \\
\hline
\end{tabularx}

\vspace{1mm}

\noindent
\begin{tabularx}{\linewidth}{|p{3.2cm}|X|}
\hline
\centering\textbf{ĐIỂM SỐ} & \centering\textbf{LỜI NHẬN XÉT CỦA THẦY CÔ} \tabularnewline
\hline
\rule{0pt}{1.3cm} & \tabularnewline
\hline
\end{tabularx}

\vspace{1.2mm}

\noindent\textbf{BẢNG ĐIỀN ĐÁP ÁN TRẮC NGHIỆM (Mỗi câu đúng 0,5 điểm):}
\vspace{0.8mm}

\noindent
\begingroup
\setlength{\tabcolsep}{3pt}
\renewcommand{\arraystretch}{1.35}
\begin{tabularx}{\linewidth}{|c|*{5}{>{\centering\arraybackslash}X|}}
\hline
\textbf{Câu} & \textbf{1} & \textbf{2} & \textbf{3} & \textbf{4} & \textbf{5} \\
\hline
\textbf{Đáp án} & & & & & \\
\hline
\end{tabularx}
\endgroup

\vspace{1.5mm}

\noindent\textbf{\large PHẦN I. TRẮC NGHIỆM KHÁCH QUAN (2,5 điểm)}
\vspace{1mm}

\cauTN{1}{Nội dung câu hỏi trắc nghiệm số 1 đang được biên soạn...}{Phương án A}{Phương án B}{Phương án C}{Phương án D}

\cauTN{2}{Nội dung câu hỏi trắc nghiệm số 2 đang được biên soạn...}{Phương án A}{Phương án B}{Phương án C}{Phương án D}

\cauTN{3}{Nội dung câu hỏi trắc nghiệm số 3 đang được biên soạn...}{Phương án A}{Phương án B}{Phương án C}{Phương án D}

\cauTN{4}{Nội dung câu hỏi trắc nghiệm số 4 đang được biên soạn...}{Phương án A}{Phương án B}{Phương án C}{Phương án D}

\cauTN{5}{Nội dung câu hỏi trắc nghiệm số 5 đang được biên soạn...}{Phương án A}{Phương án B}{Phương án C}{Phương án D}

\vspace{1.2mm}

\noindent\textbf{\large PHẦN II. TỰ LUẬN (7,5 điểm)}

\noindent\textbf{Câu 1 (2,5 điểm). Thực hiện phép tính (hợp lý nếu có thể):}
\vspace{0.8mm}

\noindent
\begin{tabularx}{\linewidth}{|X|X|}
\hline
\oBaiTap[6]{a)}{Nội dung ý a đang được cập nhật} & 
\oBaiTap[6]{b)}{Nội dung ý b đang được cập nhật} \\
\hline
\oBaiTap[6]{c)}{Nội dung ý c đang được cập nhật} & 
\oBaiTap[6]{d)}{Nội dung ý d đang được cập nhật} \\
\hline
\end{tabularx}

\newpage
% =============================================================================
% ========================= TOÁN 6 - TUẦN 26 - BUỔI 02: MÃ ĐỀ 103 =============
% ========================= MẶT SAU (TRANG 2/2) ================================
% =============================================================================

\noindent
\begin{tabularx}{\linewidth}{|X|c|}
\hline
\textbf{PHẦN II. TỰ LUẬN (Tiếp theo) - TOÁN 6 TUẦN 26 BUỔI 02} & \textbf{Mã đề: 103 - Trang 2/2} \\
\hline
\end{tabularx}

\vspace{2.5mm}

\noindent\textbf{\large Câu 2 (2,5 điểm). Tìm số tự nhiên $x$, biết:}
\vspace{0.8mm}

\noindent
\begin{tabularx}{\linewidth}{|X|X|}
\hline
\oBaiTap[6]{a)}{Nội dung ý a đang được cập nhật} & 
\oBaiTap[6]{b)}{Nội dung ý b đang được cập nhật} \\
\hline
\oBaiTap[6]{c)}{Nội dung ý c đang được cập nhật} & 
\oBaiTap[6]{d)}{Nội dung ý d đang được cập nhật} \\
\hline
\end{tabularx}

\vspace{2.5mm}

\noindent\textbf{\large Câu 3 (2,5 điểm). Bài toán thực tế:}
\vspace{0.8mm}

\noindent
\begin{tabularx}{\linewidth}{|X|}
\hline
\vspace{0.5mm}
Nội dung bài toán thực tế đang được cập nhật. Học sinh tóm tắt đề và trình bày bài giải chi tiết bên dưới.\par
\textbf{a) (1,25 điểm)} Câu hỏi ý a...\par
\textbf{b) (1,25 điểm)} Câu hỏi ý b...\par
\vspace{0.4mm}
\noindent\textit{Bài giải:}\dongke{14}
\\ \hline
\end{tabularx}

\newpage
% =============================================================================
% ========================= TOÁN 6 - TUẦN 26 - BUỔI 02: MÃ ĐỀ 104 =============
% ========================= MẶT TRƯỚC (TRANG 1/2) ==============================
% =============================================================================

\noindent
\begin{tabularx}{\linewidth}{|X|c|}
\hline
\textbf{PHIẾU KIỂM TRA ĐỊNH KỲ TOÁN 6 - TUẦN 26 - BUỔI 02} & \textbf{Mã đề: 104 - Trang 1/2} \\
\textit{Chủ đề: Bài 35: Trung điểm của đoạn thẳng} & \textit{Thời gian: 45 phút} \\
\hline
\end{tabularx}

\vspace{1mm}

\noindent
\begin{tabularx}{\linewidth}{|X|}
\hline
\rule[-3.5mm]{0pt}{10mm}\textbf{Họ và tên học sinh:} \hfill \textbf{Mã HS:} \makebox[3.5cm]{} \\
\hline
\end{tabularx}

\vspace{1mm}

\noindent
\begin{tabularx}{\linewidth}{|p{3.2cm}|X|}
\hline
\centering\textbf{ĐIỂM SỐ} & \centering\textbf{LỜI NHẬN XÉT CỦA THẦY CÔ} \tabularnewline
\hline
\rule{0pt}{1.3cm} & \tabularnewline
\hline
\end{tabularx}

\vspace{1.2mm}

\noindent\textbf{BẢNG ĐIỀN ĐÁP ÁN TRẮC NGHIỆM (Mỗi câu đúng 0,5 điểm):}
\vspace{0.8mm}

\noindent
\begingroup
\setlength{\tabcolsep}{3pt}
\renewcommand{\arraystretch}{1.35}
\begin{tabularx}{\linewidth}{|c|*{5}{>{\centering\arraybackslash}X|}}
\hline
\textbf{Câu} & \textbf{1} & \textbf{2} & \textbf{3} & \textbf{4} & \textbf{5} \\
\hline
\textbf{Đáp án} & & & & & \\
\hline
\end{tabularx}
\endgroup

\vspace{1.5mm}

\noindent\textbf{\large PHẦN I. TRẮC NGHIỆM KHÁCH QUAN (2,5 điểm)}
\vspace{1mm}

\cauTN{1}{Nội dung câu hỏi trắc nghiệm số 1 đang được biên soạn...}{Phương án A}{Phương án B}{Phương án C}{Phương án D}

\cauTN{2}{Nội dung câu hỏi trắc nghiệm số 2 đang được biên soạn...}{Phương án A}{Phương án B}{Phương án C}{Phương án D}

\cauTN{3}{Nội dung câu hỏi trắc nghiệm số 3 đang được biên soạn...}{Phương án A}{Phương án B}{Phương án C}{Phương án D}

\cauTN{4}{Nội dung câu hỏi trắc nghiệm số 4 đang được biên soạn...}{Phương án A}{Phương án B}{Phương án C}{Phương án D}

\cauTN{5}{Nội dung câu hỏi trắc nghiệm số 5 đang được biên soạn...}{Phương án A}{Phương án B}{Phương án C}{Phương án D}

\vspace{1.2mm}

\noindent\textbf{\large PHẦN II. TỰ LUẬN (7,5 điểm)}

\noindent\textbf{Câu 1 (2,5 điểm). Thực hiện phép tính (hợp lý nếu có thể):}
\vspace{0.8mm}

\noindent
\begin{tabularx}{\linewidth}{|X|X|}
\hline
\oBaiTap[6]{a)}{Nội dung ý a đang được cập nhật} & 
\oBaiTap[6]{b)}{Nội dung ý b đang được cập nhật} \\
\hline
\oBaiTap[6]{c)}{Nội dung ý c đang được cập nhật} & 
\oBaiTap[6]{d)}{Nội dung ý d đang được cập nhật} \\
\hline
\end{tabularx}

\newpage
% =============================================================================
% ========================= TOÁN 6 - TUẦN 26 - BUỔI 02: MÃ ĐỀ 104 =============
% ========================= MẶT SAU (TRANG 2/2) ================================
% =============================================================================

\noindent
\begin{tabularx}{\linewidth}{|X|c|}
\hline
\textbf{PHẦN II. TỰ LUẬN (Tiếp theo) - TOÁN 6 TUẦN 26 BUỔI 02} & \textbf{Mã đề: 104 - Trang 2/2} \\
\hline
\end{tabularx}

\vspace{2.5mm}

\noindent\textbf{\large Câu 2 (2,5 điểm). Tìm số tự nhiên $x$, biết:}
\vspace{0.8mm}

\noindent
\begin{tabularx}{\linewidth}{|X|X|}
\hline
\oBaiTap[6]{a)}{Nội dung ý a đang được cập nhật} & 
\oBaiTap[6]{b)}{Nội dung ý b đang được cập nhật} \\
\hline
\oBaiTap[6]{c)}{Nội dung ý c đang được cập nhật} & 
\oBaiTap[6]{d)}{Nội dung ý d đang được cập nhật} \\
\hline
\end{tabularx}

\vspace{2.5mm}

\noindent\textbf{\large Câu 3 (2,5 điểm). Bài toán thực tế:}
\vspace{0.8mm}

\noindent
\begin{tabularx}{\linewidth}{|X|}
\hline
\vspace{0.5mm}
Nội dung bài toán thực tế đang được cập nhật. Học sinh tóm tắt đề và trình bày bài giải chi tiết bên dưới.\par
\textbf{a) (1,25 điểm)} Câu hỏi ý a...\par
\textbf{b) (1,25 điểm)} Câu hỏi ý b...\par
\vspace{0.4mm}
\noindent\textit{Bài giải:}\dongke{14}
\\ \hline
\end{tabularx}

\end{document}
